% !PRINTED 290
\weberpage{290}
%<<CONT
oder
\[
\begin{aligned}
p\alpha' + c'\beta' &= p\alpha, & \beta' &= p\beta\\
p\gamma' + c'\delta' &= c\alpha+\gamma, & \delta' &= c\beta+\delta.
\end{aligned}
\]
Hieraus ergibt sich nach der dritten Gleichung: $c'\delta' \equiv
c\alpha+\gamma$, und mit Hilfe der vierten:
\[
c' \equiv \frac{c\alpha+\gamma}{c\beta+\delta} \quad (\mathrm{mod}\ p).
\]

Ist $c\beta+\delta \equiv 0$, so ergibt sich $c' \equiv \infty$, und ist
$c \equiv \infty$, so folgt $c' \equiv \alpha/\beta$.

Bei der Anwendung auf die Bestimmung von $c$, $c'$ kann man
die $\alpha$, $\beta$, $\gamma$, $\delta$ durch kongruente Zahlen $a$, $b$, $c$, $d$
ersetzen, deren Determinante $ad-bc$ quadratischer Rest von $c$
ist, % [sic] Gedruckt ist c; mathematisch ist p zu erwarten. und es läßt sich auch zeigen, daß man auf diese Weise jede
Substitution der Gruppe $\mathfrak{L}_0$ erhalten kann. Jede Funktion von
$v_c$ also, die durch die Substitutionen der Gruppe $\mathfrak{L}_0$ ungeändert
bleibt, bleibt daher auch ungeändert, wenn auf $\omega$ eine lineare
Substitution angewandt wird, und ist folglich eine rationale
Funktion von $j(\omega)$. Wie aber die Koeffizienten in dieser Funktion
beschaffen sind, darüber lehrt uns diese Betrachtung nichts, und
es ist daher $\mathfrak{L}_0$ nur als die Monodromiegruppe der Invariantengleichung
erkannt.

\begin{center}
\S\,79. \gesperrt{Untersuchung der Gruppe $\mathfrak{L}_0$}\textsuperscript{1}).
\end{center}

Wir haben im 10. Abschnitt des II. Bandes die Kongruenzgruppe
 und ihre Teiler ganz allgemein untersucht. Wir führen
hier diese Untersuchung, soweit sie auf das Transformationsproblem
 Bezug hat, in spezieller Form noch einmal durch.

Die in $\mathfrak{L}_0$ enthaltenen Substitutionen
\weq{1}{\begin{pmatrix}a, & b \\ c, & \partial\end{pmatrix} = \left(z, \frac{-c+az}{\partial-bz}\right),}
in denen
\weq{2}{\varDelta = a\partial - bc}
quadratischer Rest von $p$ ist, können, wie schon oben bemerkt,
auf die Form gebracht werden:
\weq{3}{\begin{pmatrix}\alpha, & \beta \\ \gamma, & \delta\end{pmatrix} = \left(z, \frac{-\gamma+\alpha z}{\delta-\beta z}\right),}
%CONT>>

\begingroup
\renewcommand{\thefootnote}{1)}
\footnotetext[1]{Über die Gruppe der linearen Substitutionen $\mathfrak{L}_0$ ist zu
vergleichen: \gesperrt{Galois}, Liouvilles Journal, Bd.\,XI. \gesperrt{Serret},
Algèbre supérieure, Section IV, Chapitre IV. C. \gesperrt{Jordan}, Traité des
Substitutions. \gesperrt{Gierster}, Gruppe der Modulargleichungen.
Mathematische Annalen, Bd.\,18.}
\endgroup
